% Standalone includable section. Requires amsmath, amssymb and a theorem environment.
\section{Adaptive shifts with a persistent reference distribution}
\label{sec:reference-shifts}

A deterministic choice among minimum circuits also has a polynomial
all-time mistake guarantee under a useful class of non-i.i.d. processes.
The condition is that the next-input distribution may adapt to the whole
past, but remains comparable to one fixed reference distribution.  In
particular, it continues to sample every reference region often enough.
The reference need not be uniform.

Let $\mathcal H$ be a finite class of $N$ Boolean hypotheses containing a
fixed target $h_\star$.  Write $\mathcal F_t$ for the information available
after round $t$, and suppose the learner chooses an
$\mathcal F_{t-1}$-measurable $h_t\in\mathcal H$ that is consistent with
all preceding examples.  The choice may depend arbitrarily on the past,
including learner randomness included in the filtration.  Its circuit
must be chosen before seeing the current input.  Labels are always
$h_\star(X_t)$.  Denote the conditional law of $X_t$ given $\mathcal F_{t-1}$ by $P_t$.
The comparison below is with respect to this same filtration, including
any learner randomness revealed before the current input.

\begin{proposition}[Adaptive reference comparability]
\label{prop:adaptive-reference}
Suppose that, almost surely, for every round $t$ and measurable set $A$,
\[
 a\mu(A)\le P_t(A)\le b\mu(A),
 \qquad 0<a\le1\le b<\infty,
\]
where $\mu$ is a fixed probability distribution.  Suppose also that every
positive disagreement probability
$\mu\{h\ne h_\star\}$, for $h\in\mathcal H$, is at least $0<\gamma\le1$.
Then the total number $M_\infty$ of mistakes satisfies
\[
 \mathbb E M_\infty
 \le b\left[1+\frac{\log N+1}{a}
                    \left(1+\log\frac1\gamma\right)\right].
\]
In particular, $M_\infty<\infty$ almost surely.
\end{proposition}

\begin{proof}
For $h\in\mathcal H$, put $p(h)=\mu\{h\ne h_\star\}$ and let
$\mathcal V_n$ be the hypotheses consistent with the first $n$ examples.
The conditional lower bound implies
\[
 \Pr(h\in\mathcal V_n)\le(1-ap(h))^n\le e^{-anp(h)}.
\]
Set $R_n=\max_{h\in\mathcal V_n}p(h)$.  By a union bound,
\[
 \Pr(R_n>u)\le\min\{1,Ne^{-anu}\}, \qquad u>0.
\]
For $0<u<\gamma$, the same bound holds with $u$ replaced by $\gamma$,
because each positive $p(h)$ is at least $\gamma$.
For any $q>0$, monotonicity gives
\[
 \sum_{n=0}^{\infty}\min\{1,Ne^{-anq}\}
 \le 1+\int_0^\infty\min\{1,Ne^{-atq}\}\,dt
 =1+\frac{\log N+1}{aq}.
\]
Integrating the tail probability over $0<u<1$, splitting at $\gamma$,
and using Tonelli's theorem therefore yields
\[
 \mathbb E\sum_{n=0}^{\infty}R_n
 \le 1+\frac{\log N+1}{a}
                     \left(1+\log\frac1\gamma\right).
\]
Conditional on $\mathcal F_{t-1}$, the learner's mistake probability is
at most $b\,p(h_t)\le bR_{t-1}$.  Summation proves the claim.
\end{proof}

For the uniform reference on $\{0,1\}^d$, one may take
$\gamma=2^{-d}$.  If hypotheses have binary descriptions of at most $L$
bits, then $N<2^{L+1}$, giving
\[
 \mathbb E M_\infty
 =O\!\left(\frac ba(L+1)(d+1)\right).
\]
This applies to an exact minimum-description consistent circuit, since a
realizing target of description length $L$ keeps every selected circuit
within that class.  For minimum gate count $s$, standard circuit counting
gives $\log N=O((s+1)\log(d+s+2))$.  No efficient algorithm for finding the exact minimum
is asserted.  For a nonuniform finite reference, $\gamma$ may instead be
bounded below by its smallest positive atom; polynomial bounds hold when
$\log(1/\gamma)$ and $b/a$ are polynomial in the problem parameters.

For example, a process of the form
\[
 P_t=(1-\eta)\mu+\eta Q_t, \qquad 0\le\eta<1,
 \qquad Q_t(A)\le B\mu(A),
\]
may choose $Q_t$ adaptively.  It satisfies the proposition with
$a=1-\eta$ and $b=1-\eta+\eta B$.  It need never become i.i.d. or converge
to a limiting distribution.

\paragraph{Summably many departures from the comparison condition.}
The same conclusion has a robust extension.  Suppose there are
nonnegative $\mathcal F_{t-1}$-measurable $\varepsilon_t$ such that
\[
 a\mu(A)-\varepsilon_t\le P_t(A)
 \le b\mu(A)+\varepsilon_t,
 \qquad \sum_{t\ge1}\varepsilon_t\le V\quad\hbox{almost surely}.
\]
Then
\[
 \mathbb E M_\infty
 \le V+b\left[1+\frac{\log N+V+1}{a}
                         \left(1+\log\frac1\gamma\right)\right].
\]
Indeed, for each fixed $h$, the process
\[
 \mathbf1\{h\in\mathcal V_n\}
 \exp\!\left(anp(h)-\sum_{t=1}^n\varepsilon_t\right)
\]
is a nonnegative supermartingale.  Thus
$\Pr(h\in\mathcal V_n)\le e^{V-anp(h)}$.
The preceding proof applies with $N$ replaced by $Ne^V$, while the
conditional mistake probability is at most $bR_{t-1}+\varepsilon_t$.
In particular, with $a=b=1$ this covers
$\sum_t\|P_t-\mu\|_{\mathrm{TV}}\le V$.
Here total variation is $\sup_A|P(A)-Q(A)|$.
This assumption controls cumulative distance from a fixed reference; it
is different from bounding only the successive changes
$\sum_t\|P_{t+1}-P_t\|_{\mathrm{TV}}$.


% Includable exact-minimum-circuit obstruction for upper-only smoothness.
\subsection{Tiny persistent shifts can defeat deterministic minimum circuits}
\label{sec:smooth-lower}

An upper density bound alone does not suffice even for exact minimum
circuits.  The failure can persist when every conditional input law is
exponentially close to uniform and, for a fixed ranking of circuits, the
sum of successive distribution changes is exponentially small.
This construction uses repeated inputs.

\begin{proposition}[Upper smoothness does not ensure few mistakes]
\label{prop:upper-smooth-lower}
Use a fixed finite fan-in-two gate basis containing AND, and count gates.
Let $d\ge4$ be divisible by four, let $\mu$ be uniform on
$\{0,1\}^d$, and put $\delta=e^{-d/8}$.
For every deterministic, history-only exact-minimum consistent circuit
rule, there is an adaptive input process with target
$f(x)=\bigwedge_{j=1}^d x_j$, of minimum size $d-1$, such that
\[
 P_t(A)\le\frac{\mu(A)}{1-\delta}\quad\hbox{for every set }A,
 \qquad \|P_t-\mu\|_{\rm TV}\le\delta
 \quad\hbox{for every round }t,
\]
yet the learner makes at least $2^{d/2}$ mistakes almost surely.
If the rule selects the first consistent minimum circuit in a fixed
ranking, the same process also satisfies
\[
 \sum_{t\ge1}\|P_{t+1}-P_t\|_{\rm TV}
 \le\frac{\delta}{1-\delta}\quad\hbox{almost surely}.
\]
There is a deterministic horizon $T$ with $\log T=O(d)$ at which
$\mathbb E M_T\ge2^{d/2-1}$.
\end{proposition}

\begin{proof}
Put $k=3d/4$, let $R=\{x:|x|<k\}$, and write
$r=|R|$, $m=\binom dk$, and $N=2^d$.
First draw independently and uniformly from $B=R\cup\{1^d\}$ until
every point of $B$ has appeared.  This stopping time is finite almost
surely.  These examples force any consistent function to depend on at
least $k$ coordinates: a function depending on a set $S$ of fewer than
$k$ coordinates cannot distinguish $1^d$ from $1_S\in R$.
A circuit with $s$ binary gates depends on at most $s+1$ coordinates.
Every $k$-variable conjunction is consistent and has $k-1$ gates, so
minimum size is exactly $k-1$.

Moreover, every minimum circuit computes a $k$-variable conjunction.
It depends on exactly $k$ coordinates, say $S$, and is one on their
all-ones assignment.  Every other assignment to $S$ occurs on a point
of weight below $k$ and is labeled zero.  Its function is therefore
precisely $\bigwedge_{j\in S}x_j$.

While any $k$-conjunction remains consistent, let $S_t$ be the support of
the learner's current minimum circuit and choose $P_t$ uniform on
$B\cup\{1_{S_t}\}$.  The current circuit is determined by the past, so
this defines a valid adaptive input process.
A draw from $B$ causes no mistake and eliminates no remaining
$k$-conjunction.  A draw of $1_{S_t}$ has conditional probability
$1/(r+2)$, causes a mistake, and eliminates exactly that conjunction:
two distinct $k$-subsets are not contained in one another.  Therefore
minimum size remains $k-1$ until all $m$ conjunctions have been
eliminated.  Those $m$ errors occur almost surely.  Thereafter retain the
final distribution, whose entire support has already been observed.

Hoeffding's inequality gives
\[
 1-r/N=\Pr\{\operatorname{Bin}(d,1/2)\ge3d/4\}
 \le e^{-d/8}=\delta.
\]
Every conditional law is uniform on a set of size $r+1$ or $r+2$.
Its density relative to $\mu$ is at most $N/(r+1)\le(1-\delta)^{-1}$,
and its total variation distance from $\mu$ is at most
$1-(r+1)/N\le\delta$, including the initial and final stages.
Under a fixed ranking, repeated points of $B$ leave the selected circuit
unchanged.  There are $m$ hard phases.  The initial switch and the
$m-1$ switches between their laws each have total variation $1/(r+2)$;
there is no final switch.  Thus total successive variation is
$m/(r+2)\le\delta/(1-\delta)$, using $m\le N-r\le N\delta$.

Finally, $m=\binom{d}{d/4}\ge(d/(d/4))^{d/4}=2^{d/2}$.
This construction permits repeated inputs; the earlier exponential
lower bounds use a permutation with no repeats. The expected time through the initial stage and all hard phases is
\[
 (r+1)H_{r+1}+m(r+2),
\]
where $H_n$ is the $n$th harmonic number.  Taking $T$ to be the ceiling of
twice this expectation gives $\mathbb E M_T\ge m/2$ by Markov's
inequality, and $\log T=O(d)$.  This horizon is exponential in $d$;
only its logarithm is linear in $d$.
\end{proof}
